
\section*{Problem 768 (Round 2)}

\subsection*{1. ROUND-2 OBJECTIVE}

Path \textbf{(C)}: isolate a genuine quantitative obstruction and strengthen the Round-1 result from mere density $0$ to an explicit subexponential upper bound.
Round 1 only proved that
\[
\frac{|A\cap[1,x]|}{x}\to 0.
\]
The exact Erd\H{o}s asymptotic
\[
\frac{|A\cap[1,x]|}{x}=\exp\!\bigl(-(c+o(1))\sqrt{\log x}\,\log\log x\bigr)
\]
remains out of reach, so the most promising Round-2 target is to close the first major gap: obtain a nontrivial upper bound on the natural scale
\[
\exp\!\bigl(-c\sqrt{\log x\,\log\log x}\bigr).
\]

\subsection*{2. Round-1 FOUNDATION USED}

I use exactly two Round-1 ingredients.

\begin{enumerate}
\item The definition of
\[
A=\Bigl\{n\ge 1:\ \forall p\mid n \text{ prime},\ \exists d\mid n,\ d>1,\ d\equiv 1 \pmod p\Bigr\}.
\]
\item The Round-1 rough/smooth decomposition idea: if $p=P(n)$ is the largest prime factor of $n\in A$, then there is a divisor of the form $1+ap$ dividing $n$, and this can be counted by summing over $(p,a)$.
\end{enumerate}

I do \emph{not} reprove the Round-1 density-$0$ theorem; I strengthen it quantitatively.

\subsection*{3. NEW INSIGHT / TOOL (ROUND 2)}

The new step is to let the smoothness threshold depend on $x$ and optimize it.

Write $\Psi(x,y)$ for the number of $y$-smooth integers $\le x$.
Then the Round-1 argument actually gives, for every $2\le y\le x$,
\[
|A\cap[1,x]|
\le
\Psi(x,y)
+
x\sum_{\substack{p>y\\ p\ \mathrm{prime}}}\ \sum_{1\le a\le x/p}\frac{1}{p(1+ap)}.
\]
Using a standard smooth-number estimate in the Canfield--Erd\H{o}s--Pomerance range, and then optimizing $y$, yields a genuine subexponential decay.

\subsection*{4. ATTACK PLAN (ROUND 2)}

The exact Round-1 gap was this: density $0$ is much too weak compared to the problem's conjectural scale.
So I now prove the following explicit theorem.

\medskip
\noindent
\textbf{Theorem.}
Let
\[
A(x)=|A\cap[1,x]|.
\]
Then, as $x\to\infty$,
\[
A(x)\le x\exp\!\Bigl(-\Bigl(\frac{1}{\sqrt 2}+o(1)\Bigr)\sqrt{\log x\,\log\log x}\Bigr).
\]

\medskip
This does \emph{not} settle Erd\H{o}s's conjectured $\sqrt{\log x}\,\log\log x$ scale, but it is a strict Round-2 advance beyond density $0$.

\subsection*{5. WORK (ROUND 2)}

Let
\[
L=\log x,\qquad L_2=\log\log x.
\]

\paragraph{Step 1: a uniform counting inequality.}
Fix $y$ with $2\le y\le x$.
Split
\[
A\cap[1,x] = A_{\le y}(x)\cup A_{>y}(x),
\]
where $A_{\le y}(x)$ consists of those $n\in A\cap[1,x]$ with largest prime factor at most $y$, and $A_{>y}(x)$ those with largest prime factor exceeding $y$.

Clearly,
\[
|A_{\le y}(x)|\le \Psi(x,y).
\]

Now let $n\in A_{>y}(x)$, and let $p=P(n)>y$ be the largest prime factor of $n$.
By the defining property of $A$, applied to this prime $p$, there exists a divisor $d\mid n$ with $d>1$ and $d\equiv 1\pmod p$.
Write
\[
d=1+ap,\qquad a\ge 1.
\]
Because $d\equiv 1\pmod p$, we have $p\nmid d$, so $p\,d\mid n$.
Hence $n$ is a multiple of
\[
p(1+ap).
\]
Therefore, for fixed $(p,a)$, the number of such $n\le x$ is at most
\[
\frac{x}{p(1+ap)}.
\]
Summing over all admissible $p$ and $a$ gives
\[
|A_{>y}(x)|
\le
x\sum_{\substack{p>y\\ p\ \mathrm{prime}}}\ \sum_{1\le a\le x/p}\frac{1}{p(1+ap)}.
\]
Since
\[
\frac{1}{p(1+ap)}\le \frac{1}{a p^2},
\]
we obtain
\[
|A_{>y}(x)|
\le
x\sum_{\substack{p>y\\ p\ \mathrm{prime}}}\frac{1}{p^2}\sum_{1\le a\le x/p}\frac{1}{a}
\le
x(1+\log x)\sum_{p>y}\frac{1}{p^2}.
\]
Finally,
\[
\sum_{p>y}\frac{1}{p^2}\le \sum_{m>y}\frac{1}{m^2}\le \frac{1}{y-1},
\]
so
\[
|A_{>y}(x)|\le \frac{x(1+\log x)}{y-1}.
\]

Combining the two pieces:
\[
A(x)\le \Psi(x,y)+\frac{x(1+\log x)}{y-1}.
\tag{5.1}
\]

\paragraph{Step 2: choose the threshold.}
Take
\[
y=\exp\!\Bigl(\frac{1}{\sqrt 2}\sqrt{L\,L_2}\Bigr).
\]
Then
\[
u:=\frac{\log x}{\log y}=\sqrt{\frac{2L}{L_2}}.
\]

A standard smooth-number theorem of Canfield--Erd\H{o}s--Pomerance (in the form quoted in Granville's survey) says that if $u\to\infty$ and $u\le y^{1-\varepsilon}$ for some fixed $\varepsilon>0$, then
\[
\Psi(x,y)=x\,u^{-u+o(u)}=x\exp\bigl(-(1+o(1))u\log u\bigr).
\]
For the present choice of $y$, we have $u\to\infty$ and, for example, $u\le y^{1/2}$ for all large $x$, because
\[
\log u=O(L_2)=o(\log y).
\]
Hence the theorem applies.

Now
\[
u\log u
=
\sqrt{\frac{2L}{L_2}}\cdot \frac{1}{2}\log\!\Bigl(\frac{2L}{L_2}\Bigr)
=
\Bigl(\frac{1}{\sqrt 2}+o(1)\Bigr)\sqrt{L\,L_2}.
\]
Therefore
\[
\Psi(x,y)
\le
x\exp\!\Bigl(-\Bigl(\frac{1}{\sqrt 2}+o(1)\Bigr)\sqrt{L\,L_2}\Bigr).
\tag{5.2}
\]

On the other hand,
\[
\frac{x(1+\log x)}{y-1}
=
x\exp\!\Bigl(-\frac{1}{\sqrt 2}\sqrt{L\,L_2}+O(L_2)\Bigr)
=
x\exp\!\Bigl(-\Bigl(\frac{1}{\sqrt 2}+o(1)\Bigr)\sqrt{L\,L_2}\Bigr),
\tag{5.3}
\]
since $L_2=o(\sqrt{L\,L_2})$.

Insert (5.2) and (5.3) into (5.1):
\[
A(x)\le x\exp\!\Bigl(-\Bigl(\frac{1}{\sqrt 2}+o(1)\Bigr)\sqrt{L\,L_2}\Bigr).
\]
This proves the theorem.

\paragraph{A simple structural check.}
The proof above is consistent with the elementary fact that every $n\in A$ satisfies
\[
P(n)\le \sqrt n.
\]
Indeed, if $p=P(n)>\sqrt n$, then $p^2>n$, so $p$ occurs only to the first power in $n$; any divisor $d\equiv 1\pmod p$ must be coprime to $p$, hence $d\mid n/p<p$, impossible because $d>1$ and $d\equiv 1\pmod p$ force $d\ge p+1$.

\subsection*{6. ADVERSARIAL VERIFICATION}

There are three places where the new argument could fail.

\begin{enumerate}
\item \textbf{Double counting of $n$.}  
That is harmless: the bound for $A_{>y}(x)$ is an upper bound, so overcounting is allowed.

\item \textbf{Range of the smooth-number theorem.}  
For my choice
\[
y=\exp\!\Bigl(\frac{1}{\sqrt 2}\sqrt{L\,L_2}\Bigr),
\]
we have $u=\sqrt{2L/L_2}\to\infty$ and
\[
\log u=O(L_2)=o(\log y),
\]
hence $u\le y^{1/2}$ for large $x$.
So the Canfield--Erd\H{o}s--Pomerance estimate applies.

\item \textbf{Possible loss in the rough part.}  
The rough bound is extremely crude, but all I need is the exponent scale.
The factor $1+\log x$ only contributes $O(L_2)$ in the exponent, which is negligible compared with $\sqrt{L\,L_2}$.
\end{enumerate}

So the Round-2 theorem is gap-free.

\subsection*{7. FINAL}

\textbf{UNRESOLVED (BUT STRICTLY ADVANCED).}

Round 2 strengthens the Round-1 conclusion from
\[
A(x)=o(x)
\]
to the explicit subexponential estimate
\[
A(x)\le x\exp\!\Bigl(-\Bigl(\frac{1}{\sqrt 2}+o(1)\Bigr)\sqrt{\log x\,\log\log x}\Bigr).
\]
This is still weaker than the upper bound quoted on the problem page, and very far from the conjectured
\[
\exp\!\bigl(-(c+o(1))\sqrt{\log x}\,\log\log x\bigr)
\]
law, but it is a real quantitative advance beyond Round 1.

\subsection*{8. COMPLETION ESTIMATE (MANDATORY)}

\textbf{COMPLETION: 38\%}

\subsection*{9. REFERENCES}

\begin{itemize}
\item T. F. Bloom, \emph{Erd\H{o}s Problem \#768}, ErdosProblems.com, accessed 2026-03-07.
\item A. Granville, \emph{Smooth numbers: computational number theory and beyond}; used for the quoted Canfield--Erd\H{o}s--Pomerance smooth-number estimate.
\end{itemize}

